\begin{mistake}{2026-04-19}{06}

\begin{problemcontent}
设 \(A = \begin{bmatrix} 1 & 1 \\ 0 & 1 \end{bmatrix}\)，则和矩阵 \(A\) 可交换的矩阵是 \underline{\hspace{3cm}}。
\end{problemcontent}

\begin{solutioncontent}
设所求矩阵为二阶方阵 \(X = \begin{bmatrix} x_1 & x_2 \\ x_3 & x_4 \end{bmatrix}\)。

由矩阵可交换的定义 \(AX = XA\)，有：
\[
\begin{bmatrix} 1 & 1 \\ 0 & 1 \end{bmatrix} \begin{bmatrix} x_1 & x_2 \\ x_3 & x_4 \end{bmatrix}
= \begin{bmatrix} x_1 & x_2 \\ x_3 & x_4 \end{bmatrix} \begin{bmatrix} 1 & 1 \\ 0 & 1 \end{bmatrix}
\]

展开两边的矩阵乘法：
\[
\begin{bmatrix} x_1 + x_3 & x_2 + x_4 \\ x_3 & x_4 \end{bmatrix}
= \begin{bmatrix} x_1 & x_1 + x_2 \\ x_3 & x_3 + x_4 \end{bmatrix}
\]

根据矩阵相等的条件，对应元素必然相等，得到线性方程组：
\[
\begin{cases}
x_1 + x_3 = x_1 \\
x_2 + x_4 = x_1 + x_2 \\
x_3 = x_3 \\
x_4 = x_3 + x_4
\end{cases}
\]

化简解得 \(x_3 = 0\)，且 \(x_4 = x_1\)。

令 \(x_1 = a\)，\(x_2 = b\)（其中 \(a, b\) 为任意常数），故和矩阵 \(A\) 可交换的矩阵为 \(\begin{bmatrix} a & b \\ 0 & a \end{bmatrix}\)。
\end{solutioncontent}

\mistakeknowledge{
\begin{itemize}
 \item \textbf{核心公式/定理}：矩阵可交换的定义 \(AX = XA\) 以及矩阵相乘法则。
 \item \textbf{易错点}：展开 \(AX\) 与 \(XA\) 时前后乘顺序弄反或元素错位；最终结果漏写“\(a, b\) 为任意常数”。
 \item \textbf{关联知识}：常考结论：矩阵 \(A\) 的任何多项式 \(f(A)\) 均与 \(A\) 可交换。
\end{itemize}}

\end{mistake}
